%% Chapter 9, Section 9.2 — Bootstrap Particle Filter
%% A First Course in Monte Carlo Methods
%% Sanz-Alonso & Al-Ghattas

\section{Bootstrap Particle Filter}
\label{sec:9.2}

In the bootstrap particle filter, we propagate particles through the dynamics model, and
then apply Bayes' formula to weight each particle according to the likelihood function
implied by the observation model. At each time-step, we resample new particles to avoid
weight degeneracy.

\begin{algorithm}[H]
  \caption{Bootstrap Particle Filter}
  \label{alg:9.1}
  \begin{algorithmic}[1]
    \Require Initial distribution $f_0^N = f_0$, observations $Y_{1:J}$, sample size $N$.
    \State For $j = 0, \ldots, J-1$ do for $1 \leq n \leq N$:
    \State Draw $X_j^{(n)} \iid f_j^N$.
    \State Propagate each sample through the dynamics model:
      \[
        \widehat{X}_{j+1}^{(n)} = a\!\left(X_j^{(n)}\right) + \xi_{j+1}^{(n)},
        \qquad
        \xi_{j+1}^{(n)} \iid \Normal(0, \Sigma).
      \]
    \State Weight each sample with the likelihood defined by the observation model:
      \[
        \tilde{w}_{j+1}^{(n)} = \exp\!\left(
          -\frac{1}{2}\Bigl|Y_{j+1} - b\!\left(\widehat{X}_{j+1}^{(n)}\right)\Bigr|_\Gamma^2
        \right).
      \]
    \State Normalize the weights:
      \[
        w_{j+1}^{(n)} = \frac{\tilde{w}_{j+1}^{(n)}}{\displaystyle\sum_{n=1}^{N} \tilde{w}_{j+1}^{(n)}}.
      \]
    \State Set
      \[
        f_{j+1}^N(x) = \sum_{n=1}^{N} w_{j+1}^{(n)} \delta\!\left(x - \widehat{X}_{j+1}^{(n)}\right).
      \]
    \Ensure Approximation $f_J^N$ to the filtering distribution $f_J$.
  \end{algorithmic}
\end{algorithm}

\medskip
\noindent\fbox{\parbox{\dimexpr\textwidth-2\fboxsep-2\fboxrule\relax}{%
\textbf{Pros and Cons}: The bootstrap particle filter is a sequential Monte Carlo
algorithm to approximate a time-evolving sequence of target distributions. It can be
applied in general hidden Markov models and it is provably convergent in the large $N$
limit. However, in high-dimensional settings the weights typically have a large variance,
leading to a small effective sample size and poor performance.
}}

\bigskip

Our goal now is to show that, for large $N$, $f_j^N$ is close to $f_j$. To perform the
analysis we first introduce three operators acting on probability distributions.

\subsection{Prediction, Analysis, and Sampling Operators}
\label{sec:9.2.1}

The first operator we introduce, which we call the prediction operator, describes how
p.d.f.s propagate through the Markov kernel specified by the dynamics model.

\begin{definition}[Prediction Operator]
  \label{def:9.1}
  The prediction operator $\mathcal{P}$ acting on a p.d.f.\ $\pi$ is given by
  \[
    (\mathcal{P}\pi)(x) = \int_{\R^d} p(\tilde{x}, x)\,\pi(\tilde{x})\,d\tilde{x},
  \]
  where $p$ is the Markov kernel associated with the dynamics model, namely
  \begin{equation}
    \label{eq:9.1}
    p(x, \tilde{x}) = \frac{1}{\sqrt{(2\pi)^d \det \Sigma}}\;
      \exp\!\left(-\frac{1}{2}\bigl|\tilde{x} - a(x)\bigr|_\Sigma^2\right).
  \end{equation}
\end{definition}

The second operator we introduce, which we call the analysis operator, takes as input a
p.d.f.\ $\pi$ and produces as output the posterior p.d.f.\ obtained by taking $\pi$ as
the prior and using the likelihood function specified by the observation model.

\begin{definition}[Analysis Operator]
  \label{def:9.2}
  The analysis operator $\mathcal{A}_j$ acting on a p.d.f.\ $\pi$ is given by
  \[
    (\mathcal{A}_j \pi)(x) \propto
      \exp\!\left(-\frac{1}{2}\bigl|Y_{j+1} - b(x)\bigr|_\Gamma^2\right)\pi(x).
  \]
\end{definition}

Denoting by $\widehat{f}_{j+1} := \mathcal{P} f_j$ the distribution of $X_{j+1} |
Y_{1:j}$, we thus have that $f_{j+1} = \mathcal{A}_j \widehat{f}_{j+1}$. Consequently,
we can write the filtering step as
\begin{equation}
  \label{eq:9.2}
  f_{j+1} = \mathcal{A}_j \mathcal{P} f_j.
\end{equation}
The operator $\mathcal{P}$ does not depend on $j$ because our dynamics model is
time-homogeneous. However, the operator $\mathcal{A}_j$ depends on the observed data
$Y_{j+1}$, and thus on the time index $j$.

\begin{remark}
  \label{rem:9.3}
  The operator $\mathcal{A}_j$ acting on a distribution of the form
  \[
    \pi(x) = \frac{1}{N} \sum \delta\!\left(x - X^{(n)}\right)
  \]
  gives
  \[
    (\mathcal{A}_j \pi)(x) = \sum_{n=1}^{N} w^{(n)} \delta\!\left(x - X^{(n)}\right),
  \]
  where
  \[
    \tilde{w}^{(n)} = \exp\!\left(-\frac{1}{2}\bigl|Y_{j+1} - b(x)\bigr|_\Gamma^2\right)
  \]
  and $w^{(n)}$ are the normalized weights.
\end{remark}

The third operator we introduce, which we call the sampling operator, takes as input a
p.d.f.\ $\pi$ and outputs the empirical p.d.f.\ obtained by drawing $N$ samples from
$\pi$.

\begin{definition}[Sampling Operator]
  \label{def:9.4}
  The sampling operator $S^N$ acts on a distribution $\pi$ by drawing $N$ samples from
  $\pi$ and producing a Dirac approximation
  \[
    (S^N \pi)(x) = \frac{1}{N} \sum_{n=1}^{N} \delta\!\left(x - X^{(n)}\right),
    \qquad X^{(n)} \iid \pi.
  \]
\end{definition}

Notice that $S^N \pi$ is a random probability distribution due to sampling from $\pi$.
With the three operators we just defined, we can succinctly write the bootstrap particle
filter as
\begin{equation}
  \label{eq:9.3}
  f_{j+1}^N = \mathcal{A}_j S^N \mathcal{P} f_j^N.
\end{equation}
In order to reconcile this way of writing the bootstrap particle filter with the
algorithmic implementation in Algorithm~\ref{alg:9.1}, it is important to note that doing
prediction and then sampling is equivalent to sampling first and then doing prediction.

\subsection{Bounds for Prediction, Analysis, and Sampling}
\label{sec:9.2.2}

We now seek to obtain bounds for the prediction, analysis, and sampling operators. We will
use the following distance between random probability distributions:
\begin{equation}
  \label{eq:9.4}
  d(\pi, \pi') = \sup_{|h|_\infty \leq 1}
    \sqrt{\mathbb{E}\!\left[\left(\mathbb{E}_\pi[h] - \mathbb{E}_{\pi'}[h]\right)^2\right]}.
\end{equation}
It can be verified that \eqref{eq:9.4} indeed defines a valid distance in the space of
random probability distributions, so that in particular it satisfies the triangle
inequality. In what follows, the randomness in the distributions we consider will arise
from sampling, and the outer expectation is with respect to such randomness.

We first show that the prediction operator is a contraction in this metric.

\begin{lemma}
  \label{lem:9.5}
  It holds that
  \[
    d(\mathcal{P}\pi, \mathcal{P}\pi') \leq d(\pi, \pi').
  \]
\end{lemma}

\begin{proof}
  Let $|h|_\infty \leq 1$ and define
  \[
    h^\dagger(x) = \int_{\R^d} p(x, \tilde{x})\,h(\tilde{x})\,d\tilde{x},
  \]
  where recall that $p$ denotes the transition kernel \eqref{eq:9.1} of the dynamics
  model. Note that
  \[
    |h^\dagger(x)| \leq \int_{\R^d} p(x, \tilde{x})\,d\tilde{x} = 1
  \]
  and
  \begin{align*}
    \mathbb{E}_\pi[h^\dagger]
    &= \int_{\R^d} h^\dagger(x)\,\pi(x)\,dx \\
    &= \int_{\R^d} \left(\int_{\R^d} p(x, \tilde{x})\,h(\tilde{x})\,d\tilde{x}\right)\pi(x)\,dx \\
    &= \int_{\R^d} \left(\int_{\R^d} p(x, \tilde{x})\,\pi(x)\,dx\right)h(\tilde{x})\,d\tilde{x} \\
    &= \int_{\R^d} (\mathcal{P}\pi)(\tilde{x})\,h(\tilde{x})\,d\tilde{x},
  \end{align*}
  and so
  \[
    \mathbb{E}_\pi[h^\dagger] = \mathbb{E}_{\mathcal{P}\pi}[h].
  \]
  Finally,
  \begin{align*}
    d(\mathcal{P}\pi, \mathcal{P}\pi')
    &= \sup_{|h|_\infty \leq 1}
       \sqrt{\mathbb{E}\!\left[\left(\mathbb{E}_{\mathcal{P}\pi}[h]
         - \mathbb{E}_{\mathcal{P}\pi'}[h]\right)^2\right]} \\
    &\leq \sup_{|h^\dagger|_\infty \leq 1}
       \sqrt{\mathbb{E}\!\left[\left(\mathbb{E}_\pi[h^\dagger]
         - \mathbb{E}_{\pi'}[h^\dagger]\right)^2\right]} \\
    &= d(\pi, \pi'),
  \end{align*}
  as desired.
\end{proof}

We next establish a bound for the analysis operator under the following assumption:

\begin{assumption}
  \label{ass:9.6}
  There exists $\kappa \in (0, 1)$ such that, for all $x \in \R^d$ and
  $j \in \{0, 1, \ldots, J-1\}$,
  \[
    \kappa \leq w_{j+1}(x) := \exp\!\left(
      -\frac{1}{2}\bigl|Y_{j+1} - b(x)\bigr|_\Gamma^2
    \right) \leq \kappa^{-1}.
  \]
\end{assumption}

While this strong assumption can be significantly relaxed, it will allow us to streamline
our theory. In particular, notice that the proof of the following result is nearly
identical to that of Theorem~3.5 for autonormalized importance sampling.

\begin{lemma}
  \label{lem:9.7}
  Under Assumption~\ref{ass:9.6},
  \[
    d(\mathcal{A}_j\pi, \mathcal{A}_j\pi') \leq \frac{2}{\kappa^2}\,d(\pi, \pi').
  \]
\end{lemma}

\begin{proof}
  Let $h$ with $|h|_\infty \leq 1$, then
  \begin{align*}
    \mathbb{E}_{\mathcal{A}_\pi}[h] - \mathbb{E}_{\mathcal{A}_{\pi'}}[h]
    &= \frac{\mathbb{E}_\pi[hw]}{\mathbb{E}_\pi[w]}
     - \frac{\mathbb{E}_{\pi'}[hw]}{\mathbb{E}_{\pi'}[w]} \\
    &= \frac{\mathbb{E}_\pi[hw]}{\mathbb{E}_\pi[w]}
     - \frac{\mathbb{E}_{\pi'}[hw]}{\mathbb{E}_\pi[w]}
     + \frac{\mathbb{E}_{\pi'}[hw]}{\mathbb{E}_\pi[w]}
     - \frac{\mathbb{E}_{\pi'}[hw]}{\mathbb{E}_{\pi'}[w]} \\
    &= \frac{1}{\kappa}\!\left(
       \frac{\mathbb{E}_\pi[h\kappa w] - \mathbb{E}_{\pi'}[h\kappa w]}{\mathbb{E}_\pi[w]}
       + \frac{\mathbb{E}_{\pi'}[hw]}{\mathbb{E}_{\pi'}[w]}
         \cdot
         \frac{\mathbb{E}_{\pi'}[\kappa w] - \mathbb{E}_\pi[\kappa w]}{\mathbb{E}_\pi[w]}
      \right).
  \end{align*}
  Since $|h|_\infty \leq 1$, we have
  \[
    \left|\mathbb{E}_{\mathcal{A}_{\pi'}}[h]\right|
    = \left|\frac{\mathbb{E}_{\pi'}[hw]}{\mathbb{E}_{\pi'}[w]}\right| \leq 1.
  \]
  In addition, by Assumption~\ref{ass:9.6}, we have $\mathbb{E}_\pi[w] \geq \kappa$.
  Therefore,
  \[
    \left|\mathbb{E}_{\mathcal{A}_\pi}[h] - \mathbb{E}_{\mathcal{A}_{\pi'}}[h]\right|
    \leq \frac{1}{\kappa^2}\!\left(
      \left|\mathbb{E}_\pi[h\kappa w] - \mathbb{E}_{\pi'}[h\kappa w]\right|
      + \left|\mathbb{E}_{\pi'}[\kappa w] - \mathbb{E}_\pi[\kappa w]\right|
    \right),
  \]
  and so
  \begin{align*}
    \mathbb{E}\!\left[\left(\mathbb{E}_{\mathcal{A}_\pi}[h]
      - \mathbb{E}_{\mathcal{A}_{\pi'}}[h]\right)^2\right]
    &\leq \frac{2}{\kappa^2}\!\left(
       \mathbb{E}\!\left[\left(\mathbb{E}_\pi[h\kappa w]
         - \mathbb{E}_{\pi'}[h\kappa w]\right)^2\right]
       + \mathbb{E}\!\left[\left(\mathbb{E}_{\pi'}[\kappa w]
         - \mathbb{E}_\pi[\kappa w]\right)^2\right]
      \right) \\
    &\leq \frac{4}{\kappa^4}\,d(\pi, \pi')^2,
  \end{align*}
  where in the last line we have used that $|\kappa w| \leq 1$. Taking square roots and
  the supremum over $|h|_\infty \leq 1$ completes the proof.
\end{proof}

Lastly, the sampling operator can be bounded using Theorem~3.1 for classical Monte Carlo
integration.

\begin{lemma}
  \label{lem:9.8}
  For any distribution $\pi$,
  \[
    d(\pi, S^N \pi) \leq \frac{1}{\sqrt{N}}.
  \]
\end{lemma}

\begin{proof}
  For any $h$ with $|h|_\infty \leq 1$,
  \[
    \mathbb{E}_{S^N\pi}[h] = \frac{1}{N}\sum_{n=1}^{N} h\!\left(X^{(n)}\right),
  \]
  where $X^{(n)} \iid \pi$. Therefore, using Theorem~3.1 for classical Monte Carlo
  integration,
  \[
    \mathbb{E}\!\left[\left(\mathbb{E}_{S^N\pi}[h] - \mathbb{E}_\pi[h]\right)^2\right]
    = \mathbb{E}\!\left[\left(\calI_\pi^{\mathrm{mc}}[h] - \calI_\pi[h]\right)^2\right]
    = \frac{\mathbb{V}_\pi[h]}{N}
    \leq \frac{1}{N}.
  \]
  Taking square roots and the supremum over $|h|_\infty \leq 1$ gives the result.
\end{proof}

\subsection{Error of Bootstrap Particle Filter}
\label{sec:9.2.3}

Recall from equations \eqref{eq:9.2} and \eqref{eq:9.3} that the filtering step and the
bootstrap particle filter can be succinctly written as
\begin{align}
  f_{j+1} &= \mathcal{A}_j \mathcal{P} f_j, \label{eq:9.5} \\
  f_{j+1}^N &= \mathcal{A}_j S^N \mathcal{P} f_j^N. \label{eq:9.6}
\end{align}
With the bounds we derived for prediction, analysis, and sampling operators in
Lemmas~\ref{lem:9.5}, \ref{lem:9.7}, and \ref{lem:9.8}, the following theorem controls
the error of the bootstrap particle filter. The proof resembles the classical Lax
equivalence framework in numerical analysis, where convergence is established by ensuring
consistency and stability.

\begin{theorem}[Error of Bootstrap Particle Filter]
  \label{thm:9.9}
  Under Assumption~\ref{ass:9.6}, there exists a constant $c = c(J, \kappa)$ such that,
  for $j = 1, \ldots, J$,
  \[
    d(f_j, f_j^N) \leq \frac{c}{\sqrt{N}}.
  \]
\end{theorem}

\begin{proof}
  Let $e_j = d(f_j, f_j^N)$, then using \eqref{eq:9.5} and the triangle inequality for
  the distance $d$ defined in \eqref{eq:9.4}, we have
  \begin{align*}
    e_{j+1}
    &= d(f_{j+1}, f_{j+1}^N)
     = d(\mathcal{A}_j \mathcal{P} f_j,\, \mathcal{A}_j S^N \mathcal{P} f_j^N) \\
    &\leq d(\mathcal{A}_j \mathcal{P} f_j,\, \mathcal{A}_j \mathcal{P} f_j^N)
     + d(\mathcal{A}_j \mathcal{P} f_j^N,\, \mathcal{A}_j S^N \mathcal{P} f_j^N).
  \end{align*}
  The first term represents the growth of error over one time-iteration by propagating
  through the operator $\mathcal{A}_j \mathcal{P}$ which governs the true evolution of
  the filtering distribution, whereas the second term is a truncation error, i.e.\ the
  error incurred by the algorithm in one iteration. We then have
  \begin{align*}
    e_{j+1}
    &\leq \frac{2}{\kappa^2}\!\left(
       d(\mathcal{P} f_j^N, \mathcal{P} f_j)
       + d(\mathcal{P} f_j^N, S^N \mathcal{P} f_j^N)
      \right) \\
    &\leq \frac{2}{\kappa^2}\!\left(e_j + \frac{1}{\sqrt{N}}\right),
  \end{align*}
  where in the first inequality we used Lemma~\ref{lem:9.7} for the analysis operator and
  in the second inequality we used Lemmas~\ref{lem:9.5} and \ref{lem:9.8} for the
  prediction and sampling operators. Letting $\lambda = \frac{2}{\kappa^2}$, it then
  follows from induction that
  \[
    e_j \leq \lambda^j e_0 + \frac{\lambda}{\sqrt{N}} \cdot \frac{1 - \lambda^j}{1 - \lambda},
  \]
  where $e_0 = 0$ since $f_0^N = f_0$. Now, noting that the above expression is
  monotonically increasing in $j$ completes the proof by setting
  $c = \frac{\lambda(1 - \lambda^J)}{1 - \lambda}$.
\end{proof}
